\section{Discussion}\label{sec:discussion}

Two results of this paper are exact. The compressive sector of KRMHD is a continuum of passive slices labelled by phase velocity, for each branch and in truncated hierarchies, for any smooth Alfv\'enic fields and without collisions (theorems~\ref{thm:slice} and~\ref{thm:trunc}), so the nonlinearity moves no free energy between phase velocities. With an ideal induction equation, smooth fields and no spatial dissipation of the compressive distribution, every volume-integrated single-point statistic, including $W(m)$ and the field-aligned Hermite flux, evolves as in linear phase mixing along the initial field lines (theorem~\ref{thm:frozen}), so such dynamics produces no stochastic echoes in $W(m)$.

The two-channel model extends these results to turbulence, where they do not apply directly, by treating spatial dissipation as a sink that removes free energy from every Hermite moment at the passive-scalar cascade rate. In the pilot simulation at $\beta_i=1$ and $\chi\approx1$ the model holds to the accuracy of the measurement. The perpendicular dissipation is close to an $m$-independent sink, and $W(m)$ and the field-aligned flux are those of phase mixing at the injected field-aligned wavenumber with that sink, with no deficit at large $m$. The density moment acquires substantial structure along the field only at the perpendicular scales where the hyper-resistivity acts. In this run the relabelling of field lines that \citet{Meyrand2019} invoke happens, but only to free energy that is being removed at the same scales, and the free energy that phase mixes keeps its injected wavenumber.

Our spectrum at $\beta_i=1$ differs from that of \citet{Meyrand2019}, who found $W(m)$ somewhat steeper than $m^{-1}$ at the same $\beta_i$ and a net flux to high $m$ much smaller than the injection rate. \placeholder{compare the two set-ups: forcing (fields, band, correlation time), dissipation of $g$ ($k_\perp^4$ here, $k_\perp^8$ there), resolution in $k_\perp$ and $m$, and the flux diagnostic (summed over $k_\perp$ here, per $k_\perp$ shell there); state which differences could change $W(m)$ and which runs would decide it} Two points are already clear. A diagnostic built on the Eulerian $k_z$ reports a backward flux of about 12\% of the forward flux in our run, where the field-aligned measures show none, so splits and fluxes of the free energy should be computed along the perturbed field. And the field-aligned wavenumber of $g$ measured with $k_z$ grows with $k_\perp$ much as that of the Alfv\'enic fields does, whereas measured along the field lines it stays close to the injected value through most of the inertial range.

The model predicts that $m_*$ grows in proportion to $\beta_i$ at fixed $\chi$, which the pilot at $\beta_i=1$ cannot test because there $m_*\gg M$. \placeholder{state the outcome of the $\beta_i$ scan (section~\ref{sim:beta}): whether the break enters the hierarchy at $\beta_i=0.3$ and $0.1$ where predicted, and how this relates to the steeper spectra reported at $\beta_i=0.1$ by \citet{Kawazura2019,Kawazura2020}} The trend those authors saw, $m^{-1/2}$ at high $\beta_i$ and close to $m^{-1}$ at $\beta_i=0.1$, is the trend the model predicts without any echo, because $m_*$ falls with $\beta_i$.

The main open question is the limit of high resolution. The model needs a perpendicular cascade whose rate does not depend on the dissipation coefficients and field lines that stay frozen for the free energy that phase mixes. \citet{DrivasEyink2017} and \citet{Eyink2018} give reasons to expect that the first requires the failure of the second at small enough scales (section~\ref{tc:limits}). In our run $T=15.9\,\tau_{\rm nl}$ at $64^3$, with an inertial range of a factor of four in $k_\perp$; whether $T$ and the frozen wavenumber of $g$ survive a longer inertial range is what section~\ref{sim:resolution} must decide. \placeholder{result of the resolution study}

The results hold for $k_\perp\rho_i\to0$. At finite $k_\perp\rho_i$ the Landau coupling depends on the wavenumber, the advection depends on $v_\perp$ and the compressive fluctuations act back on the Alfv\'enic ones \citep{Chandran2026}, so the slices no longer decouple. Velocity-dependent collisions are the only terms in KRMHD that exchange free energy between slices (\ref{eq:sliceforcing}).

% Acknowledgements to be written by the author.
